% page 552 du fac-simile (image p-551 du scan)
% bloc 160.0 x 242.0 mm ; marge gauche 8.0 mm
\pgc{-12.700mm}{0.000mm}{
\l{\cel{3}544.}
\l{}
\l{\sou{spanyoleto}. Fero-peco qua uzesas por klozar od apertar fenes-}
\l{\cel{3}tro per mancho manovrebla rotace.- DFI.}
\l{}
\l{\sou{sparar}. (\sou{trans.}) Ne spensar la toto, por povar rezervar la}
\l{\cel{3}restajo. - D.}
\l{}
\l{\sou{sparadrapo.(medic.})\cel{2}Texuro di qua la averso provizesas ye}
\l{\cel{3}strato de plastro. -}
\l{}
\l{\sou{sparviero}. (\sou{zool}.) Speco di falkono quan onu dresas por chaso.}
\l{\cel{3}- L.\cel{1}\sou{falco nisus}. - DFI.}
\l{}
\l{\sou{spasmo}. (\sou{medic.}) Kontrakteso bruska di ula organi (stomako,}
\l{\cel{3}intestini, veziko, e c.) - EFIRS.}
\l{}
\l{\sou{spato}. (\sou{mineral.}) Nomo di plura substanci minerala lameloza}
\l{\cel{3}e kolor-chanjanta. - DFIRS.}
\l{}
\l{\sou{spatul}o. Instrumento de ligno, ivoro, metalo, e c., plu larja}
\l{\cel{3}e plata ye un extremajo kam kuliero, quan uzas (1) la apo-}
\l{\cel{3}tekisti por agitar la preparaji farmaciala, extensar la}
\l{\cel{3}unguenti, plastri, e c.; (2) la piktisti, por dilutar e}
\l{\cel{3}pistar la farbi ; (c) la skultisti, por modlar vaxo, ar-}
\l{\cel{3}gilo ; (d) la gisisti, por facar mulduri. - DEFIS.}
\l{}
\l{\sou{speco}. Dividuro di genero. - DEFIRS.}
\l{}
\l{\sou{specifika}. Qua indikas ta o ca speco, exkluzante omna ceteri.}
\l{\cel{3}DEFIS.}
\l{}
\l{\sou{specimeno}. Parto de ensemblo destinita a donar ideo pri la}
\l{\cel{3}cetero. - DEFS.}
\l{}
\l{\sou{spegul}o. Vitro-plako, polisita, di qua la reverso esas kovrita}
\l{\cel{3}indute ye folio de stano merkuriizita, o metalo-plako poli-}
\l{\cel{3}sita, ube onu povas vidar sua vizajo reflektita. -DEI.}
\l{}
\l{spektar. (\sou{trans.)}. Vidar (ulo) eventar avan sua regardo. -EFIS}
\l{}
\l{spektaklo. I. Reprezentajo teatrala. - II. Ensemblo quan em-}
\l{\cel{3}bracas la regardo. - EFIRS.}
\l{}
\l{spektro.\cel{3}I.Ento nevidebla qua aparas, pavorigante, terori-}
\l{\cel{3}give, lor halucino. - II. (\sou{optiko}). Imajo oblonga en qua}
\l{\cel{3}reflektesas, forme di bendi interparalela, reda, oranjea,}
\l{\cel{3}flava, verda, blua - la radii parta qui rezultas de la des-}
\l{\cel{3}kompozo di radio ek lumo sen-kolora trans prismo de vitro.}
\l{\cel{3}- (en la senco I, preferindan esas uzar \textquotedbl{}fantomo\textquotedbl{}). -}
\l{\cel{3}DEFIRS.}
\l{}
\l{spektroskopo. (\sou{fiziko}). Aparato quan onu uzas por studiar la}
\l{\cel{3}lumo-spektri. - DEFIRS.}
\l{}
\l{spektroskopio. (\sou{fiz.}) Analizo studia di la lumo-spektri.-DEFIRS}
}
